% Included by mathematical_model.tex; evidence is frozen at the dated snapshot.
\clearpage
\section{PV forecasting: experiments, formulations and first results}
\label{sec:forecast-research}

\subsection{Targets, information and the meaning of a forecast error}
Four anonymous farms have AC capacities 0.95328, 1.5246, 4.16 and 6.638 MW.
Location, some geometry and every hypothetical battery/POI parameter remain
modeling assumptions. A weather initialization, its assumed release, the target
interval, the time a target becomes observable, and the time a fitted model
becomes available are different clocks. For a decision at $k$, admissibility is
\begin{equation}
 A_v\le k,\qquad A_{\rm model}\le k,
 \qquad A_{\rm label}\le k\quad\text{for any feedback used in fitting/updating}.
\end{equation}
Historical studies can explicitly relax model availability, but must retain the
retrospective label. The six-hour weather-release convention is assumed.
An initialized forecast is not automatically available at initialization.

The reviewed dataset distinguishes retained measured AC, approved synthetic
replacement, policy-defined zero, and still-excluded/missing intervals. A
historical-GHI physical reconstruction estimates a counterfactual production
profile; it does not establish that an observed low-output episode was faulty.
The newer optimization profile fills retained daylight exclusions with its
stored historical-GHI reference and fails on unresolved daylight gaps. The
training and scoring masks are specific to each study. Scores on old raw/curated
targets and newer reviewed targets cannot be compared as a controlled ablation.

For quarter $t$, normalized target is $y_t=P_t/P_{\rm AC}$. With eligible daylight
quarters $\mathcal I_d$ on day $d$, the reviewed search uses day-balanced scores
\begin{align}
 \mathrm{nRMSE}&=100\sqrt{\frac1{|\mathcal D|}\sum_d
             \frac1{|\mathcal I_d|}\sum_{t\in\mathcal I_d}(\hat y_t-y_t)^2},\\
 \mathrm{nMAE}&=\frac{100}{|\mathcal D|}\sum_d
             \frac1{|\mathcal I_d|}\sum_{t\in\mathcal I_d}|\hat y_t-y_t|.
\end{align}
Signed bias replaces the absolute residual by $\hat y-y$. Day-balanced and
quarter-weighted summaries need not coincide. A day's measured and synthetic
subsets may both be present, so subgroup day counts need not sum to the all-target
day count. Auction views overlap and cannot be summed as independent production.

\subsection{Physical calibration and irradiance diagnostics}
The original physical-calibration workflow minimizes a bounded, normalized
day/auction loss, optionally regularized toward reference assumptions:
\begin{equation}
 \min_{\theta\in\mathcal B}
 \frac1{|\mathcal G|}\sum_{(d,a)\in\mathcal G}\frac1{|\mathcal I_{da}|}
 \sum_{t\in\mathcal I_{da}}\rho\!\left(
 \frac{f(X_{dat};\theta)-P_{dt}}{P_{\rm AC}}\right)
 +\lambda\sum_j\left(\frac{\theta_j-\theta_{0j}}{u_j-l_j}\right)^2.
\end{equation}
Candidate families fit an AC scale, Faiman coefficients and, for trackers,
ground-coverage ratio and maximum rotation; bifacial candidates additionally
state rear-irradiance assumptions. Monthly training precedes the issue, the
last seven completed days select history/transposition/regularization, and
the selected candidate is refitted on the eligible history. The original
2025 study used January warmup and 334 February--December evaluation days.
The earlier feature-tree LaTeX records held-out daylight RMSE reductions from
0.189674 to 0.076846 MW (farm 1), 0.333161 to 0.144849 (345), 0.632407 to
0.571855 (901) and 1.119012 to 1.032802 (910). These are retained historical
report values; the full original physical-calibration run directory is not
present in the inspected main-tree output inventory, so they are secondary
evidence rather than a newly reconstructed result. The final all-year parameter
fits are in-sample and were treated as prior effective metadata. Current shared
temperature defaults differ from those historical fitted thermal coefficients.

A later retrospective physical study used historical satellite GHI and the
provider-confirmed historical diffuse fraction, forecast temperature/wind from
four cycles, and a January 2025--July 2026 fitting period. August 2026 selected
the configuration. It evaluated 624 initial trials and 12 refinements; bounded
L-BFGS-B optimized thermal coefficients, temperature coefficient and scale.
The four weather views each carried weight 0.25; they were correlated views,
not four observed weather realizations. Daytime RMSE normalized by AC rating was:
\begin{center}\small
\begin{tabular}{rllrr}
\toprule Farm & Selected proxy & Selected configuration & Original & Selected\\\midrule
1 & Milan & Unchanged reference & 3.416\% & 3.416\%\\
345 & Genoa & Fixed, isotropic, Huber & 13.493\% & 8.675\%\\
901 & Trieste & Rear-enabled tracker, Huber & 8.192\% & 7.290\%\\
910 & Aosta & Tracker, isotropic, MSE & 14.027\% & 13.151\%\\
\bottomrule\end{tabular}\end{center}
These are historical-irradiance conversion scores, not operational forecast
scores. Several fitted coefficients reached bounds, structural trials were
unbalanced, and August was used for selection. The saved report
\path{docs/results/pv_realized_irradiance_2025_2026/README.md} preserves those
qualifications. In a matched August 20--31 daylight forecast comparison,
January-only quantile calibration gave new/previous-physical CRPS of
9.970/9.970, 8.901/9.692, 7.306/7.878 and 7.938/7.850 percent AC respectively.
The unchanged farm-1 physics had identical January-only scores. Its much better
earlier calibrated score used additional calibration history; it does not
identify a physical-model improvement.

The older neural-pilot study trained on 2025, validated on January--April 2026
and inspected May--August. Timing ablations added lead and cycle hour, with
80 fits across five farm/family candidates. The retained TCN choices were partly
informed by the later period, so they are research pilots. Satellite-GHI input
swaps diagnose sensitivity to irradiance error; they leave other weather and
unobserved operation uncertain. In particular, farm 910's August low-output
episodes were not automatically diagnosed as faults by this exercise.

\subsection{Completed reviewed-GHI model search}
The completed run \path{historical_ghi_all_farms_night3000_v1} contains 8,192
candidate fits: four farms, two feature sets, sixteen families and 64 sampled
settings per group. It completed 132 selected-model refits. Families cover
regularized linear/robust regression, spline/polynomial/random-feature ridge,
random/extra trees, histogram/LightGBM/CatBoost boosting, MLP, TCN,
bidirectional GRU/LSTM and Transformer. This is a sampled search, not a global
architecture optimum. Training used 2025, validation January--April 2026 and
test May--August, with the saved boundary embargo.

Training substitutes historical GHI for forecast GHI, while other weather
comes from the preceding delivery day's 00 UTC cycle. Validation/test use
forecast GHI. Hourly mean radiation is repeated over its constituent quarters;
instantaneous weather is interpolated to quarter-hour midpoints. Preprocessing
is fitted on training only. Exactly 3,000 eligible physical nighttime training
quarters per farm are selected across the year. Daylight performance controls
selection; 2026 has already been inspected and contains synthetic labels.

For a neural predictor, let $X$ be the available weather sequence and let
$f_\theta(X)_t$ be its normalized AC prediction. The core learning problem is
\begin{equation}\label{eq:pv-training}
 \min_\theta \sum_{d,t}\omega_{dt}\,
ho\bigl(f_\theta(X)_t-y_{dt}\bigr),
 \qquad \sum_{t\in d}\omega_{dt}=\frac1{|\mathcal D_{\rm train}|},
 \quad \hat P_t=P_{\rm AC}\clip(f_\theta(X)_t,0,1).
\end{equation}
The selected family/hyperparameters define MSE or Huber loss; the neural Huber
transition is 0.05 of AC capacity. AdamW trains up to 150 epochs with patience
15 on validation daylight error. The training objective and early-stopping
criterion are separate from the final test score; there is no post-hoc
calibration in this point-model search. Monthly night windows force zero output.

An MLP applies successive affine maps and nonlinearities to weather features.
A TCN combines finite dilated convolutions over the supplied forecast sequence.
A bidirectional recurrent model concatenates forward and backward hidden states
of that same issued forecast. Transformer attention has the usual form
\begin{equation}
 H=\operatorname{softmax}(QK^{\mathsf T}/\sqrt{d_k})V,
 \qquad Q=XW_Q,\ K=XW_K,\ V=XW_V,
\end{equation}
with the configured positional, masking and feed-forward transformations.
Using future \emph{forecast} covariates in a sequence is not future-target
disclosure when the complete forecast was issued by the cutoff. In contrast,
using future observed GHI in validation would be disclosure and is excluded.
Sequence boundaries and shortened inference context remain material choices.

The following validation-selected winners use extended features. Scores are
day-balanced daylight percentages of AC capacity. Test ranking was not used
to replace the validation winner.
\begingroup\small
\begin{longtable}{llrrrr}
\toprule
Farm & Model & Validation & Test & Test MAE & Test bias\\
\midrule\endhead
1 & bigru & 4.708 & 6.069 & 4.006 & -1.048\\
345 & transformer & 7.645 & 10.366 & 6.893 & -1.169\\
901 & bilstm & 10.241 & 11.231 & 6.776 & 1.137\\
910 & tcn & 8.977 & 10.992 & 6.630 & 1.071\\
\bottomrule
\end{longtable}
\endgroup

Zero and a training-only day-balanced constant are included in the saved
baseline table. The 132 saved selected models were subsequently rescored with
nighttime clipping; daylight metrics and all 8,192 validation rankings were
unchanged. The retained inference bundles are the four winners, not every
historical fitted weight file. Test errors separated by target origin are:
\begingroup\small
\begin{longtable}{lrrrrr}
\toprule
Farm & Test target & Quarters & Days & nRMSE (\%) & nMAE (\%)\\
\midrule\endhead
1 & all & 8094 & 123 & 6.069 & 4.006\\
1 & measured & 8065 & 123 & 6.067 & 4.004\\
1 & synthetic & 29 & 3 & 3.977 & 2.993\\
345 & all & 8068 & 123 & 10.366 & 6.893\\
345 & measured & 8058 & 123 & 10.364 & 6.888\\
345 & synthetic & 10 & 3 & 10.179 & 8.482\\
910 & all & 8075 & 123 & 10.992 & 6.630\\
910 & measured & 5735 & 90 & 12.045 & 7.326\\
910 & synthetic & 2340 & 41 & 8.702 & 5.615\\
901 & all & 8108 & 123 & 11.231 & 6.776\\
901 & measured & 8037 & 123 & 11.251 & 6.764\\
901 & synthetic & 71 & 7 & 6.695 & 4.613\\
\bottomrule
\end{longtable}
\endgroup

Farm 910 has 2,340 synthetic scored test quarters; its all-target 10.992\%
nRMSE cannot be interpreted as independent metered accuracy. Its measured-only
12.045\% and synthetic-only 8.702\% describe different populations, not a paired
estimate of reconstruction bias. Seven-day block bootstrap comparisons are
conditional on the chosen targets and previously inspected dates.

\subsection{Four-cycle, 48-hour PV extension in progress}
The new neural study uses cycles at 00, 06, 12 and 18 UTC and 192 quarter-hour
targets per 48-hour issue window. A physical target can occur in multiple
vintages. All copies retain the target's chronological split, and the first
six hours precede the assumed release. These early leads can be retrospective
training examples but are excluded from availability-eligible validation
selection. The proposed training objective extends~\eqref{eq:pv-training} to
issue/lead pairs $(v,\ell)$ with $t=t_v+\ell$; the effective sampling weights
are defined by the saved data preparation, not by assuming overlapping views
are independent observations. Uncertainty reporting must block by target day.

The checked dataset has, per farm, 195,264 training pairs (including 24,000 views
of 3,000 selected nighttime quarters), 47,656 validation pairs and 56,826 test
pairs. Core weather must be present for a supervised row; additional gaps use
training medians, with no missingness indicator channels in this design.
Boundary sequence-context missingness remains reported. The five neural
families, two feature sets and 64 draws make 2,560 candidate fits, followed by
40 family/feature refits and up to eight extra top-three refits.

The CPU, single-GPU, parallel-GPU and multi-GPU attempts have distinct identities.
The multi-GPU search was interrupted by GPU-memory contention; the recovery
record at this snapshot has 220/2,560 completed candidates and mode
\texttt{gpu0\_recovery}. Other saved partial attempts must not be pooled into a
completed search. The configured parallel recovery uses two fit workers on
GPU 0 and retains matching earlier records. No final four-cycle ranking or
neural-controller result is available in the captured evidence. Synthetic GPU
throughput checks are engineering measurements, not model-accuracy experiments.

\subsection{Marginal ensemble calibration and delayed online CRPS}
The completed January/August marginal study uses 49 training and 12 validation
delivery days. Forecasts from 50 weather members are converted to PV; model
selection separately ranks the member mean by RMSE, median by MAE, and calibrated
quantiles by CRPS. For quantile level $\tau$, the pinball objective is
\begin{equation}
 \min_\beta\sum_i \rho_\tau(y_i-x_i^{\mathsf T}\beta),
 \qquad \rho_\tau(r)=r\bigl(\tau-\mathbf1\{r<0\}\bigr).
\end{equation}
Quantiles are bounded/sorted as prescribed by that study, linearly interpolated
between saved levels, with constant tails beyond p05 and p95. This is a
marginal distribution; it does not identify a trajectory law.

For expert predictions $x_m$ and probability weights $w_m$, exact weighted
empirical CRPS is
\begin{equation}\label{eq:weighted-crps}
 \operatorname{CRPS}(w,x,y)=\sum_m w_m|x_m-y|
 -\frac12\sum_{m,j}w_mw_j|x_m-x_j|.
\end{equation}
Quantile-curve CRPS instead integrates twice the pinball loss over probability.
Finite-grid quadrature is not an exact empirical-ensemble score.

The delayed ML-Poly implementation uses normalized powers and
\begin{align}
 g_m&=|x_m-y|-\sum_jw_j|x_m-x_j|,&
 r_m&=\sum_jw_jg_j-g_m,\\
 R_m&\leftarrow R_m+r_m,&S_m&\leftarrow S_m+r_m^2,\\
 \eta_m&=(1+S_m)^{-1},&
 w_m^{\rm next}&=\frac{\eta_m\pos{R_m}}{\sum_j\eta_j\pos{R_j}}.
\end{align}
A zero denominator uses uniform weights. Learners are separate by farm,
auction and three-hour local block. Expert identities are sorted ranks: 50
ensemble-only experts, or 70 including a deterministic point and 19 error
quantiles. Warm starts replay eligible earlier out-of-fold feedback; cold
starts begin uniformly. Each issue freezes weights. Delayed gradients use the
weights originally issued, and feedback is released only after target
availability. No regret theorem is claimed for this modified delayed replay.

On August 20--31, the DA CRPS comparison is:
\begin{center}\small
\begin{tabular}{rrrrr}
\toprule Farm & Matched quarters & Equal ensemble & Fixed quantiles & Warm combined\\\midrule
1 & 1152 & 2.493 & 2.048 & 2.120\\
345 & 1152 & 2.774 & 2.418 & 2.704\\
901 & 1152 & 4.792 & 4.307 & 4.630\\
910 & 1101 & 2.793 & 2.288 & 2.503\\
\bottomrule\end{tabular}\end{center}
Units are percent AC; these full matched samples include night. Fixed quantiles
beat warm online combination at DA on all four farms. At IDA3, warm combination
improves fixed-quantile CRPS for farms 1 (2.318 versus 2.403) and 910 (2.204
versus 2.352), but not 345 or 901. This is mixed exploratory evidence, not a
general online-learning benefit. Candidates were selected on these validation
dates; August 20 early auctions include a retrospective source-model case.
The complete 112-row comparison is retained in the original metrics file.

The matched same-day-cycle study is a separate hindsight sensitivity: the
current day's 00 UTC cycle is unavailable at that day's DA/IDA1 cutoffs.
Calibrated CRPS changes from preceding-day to same-day were 2.05 to 1.89, 2.42
to 2.55, 4.31 to 3.86 and 2.29 to 2.09 percent AC for farms 1, 345, 901 and
910. Selection was repeated on the same 12 days, so this is neither an isolated
vintage ablation with frozen models nor evidence of a tradable DA improvement.

\subsection{Censored-distribution and quantile-network extension}
The Mayer-style extension is implemented but no completed historical run
directory was found in the inspected outputs. Its current broader configuration
requires ensemble and target coverage beyond the available August boundary.
It therefore has no new numerical result in this report. The input is the
normalized hourly PV ensemble plus a deterministic control \emph{proxy}; only
complete valid four-quarter hours enter. The proxy is not an identified
meteorological control member.

For $Z\sim N(\mu,\sigma^2)$, doubly censored PV is
\begin{equation}
 Y=\min(1,\max(0,Z)),\quad
 \Pr(Y=0)=\Phi(-\mu/\sigma),\quad
 \Pr(Y=1)=1-\Phi((1-\mu)/\sigma).
\end{equation}
Basic EMOS links location to predictors and log scale to ensemble variance,
fitting per lead by CRPS. Boosted censored regression pools leads and fits
likelihood. Linear quantile regression uses sorted members and the proxy.
The neural censored-distribution model predicts $\mu_\theta(x)$ and
$\sigma_\theta(x)=\operatorname{softplus}(h_\theta(x))+\epsilon$ and minimizes
mean censored-normal CRPS. Ten seeded models average location and scale;
this is not a ten-component mixture distribution. The analytic PyTorch score
retains native gradients and is checked against R \texttt{scoringRules}; R
\texttt{crch} and scoring run in isolated subprocesses~\cite{scoring}.

QRNN predicts 51 quantiles using a smooth asymmetric quantile loss. A Bernstein
quantile network imposes monotonicity through ordered coefficients:
\begin{equation}
 Q_\theta(\tau)=\sum_{j=0}^K a_j(x)\binom{K}{j}\tau^j(1-\tau)^{K-j},
 \qquad a_j=a_0+\sum_{h=1}^j\operatorname{softplus}(b_h).
\end{equation}
The noncrossing network uses a positive triangular cumulative construction.
Its architecture is an explicit implementation choice, not an exact paper
reproduction. Out-of-range quantiles are flagged/excluded from matched scores
rather than silently clipped in this post-processing study. The quantile
families have 100 seeded random trials per farm, inner five-day blocks for
training/stopping/selection and separate outer evaluation. No copula or MILP
scenario coupling is produced by these marginal models alone.

\section{Price forecasting across the two sibling workflows}
\label{sec:price-research}

\subsection{Feature data and availability contract}
The feature checkout builds snapshot-pinned market/fundamental/weather tables.
Its 2024 coverage report contains 8,784 native hourly NORD DA labels with no
missing quarter-hour-equivalent slots; load and total-generation forecasts also
cover the year. Solar and wind each miss 388 quarter-hour slots. The French
national balance misses four such slots at the autumn DST transition. Gas IGI
has 366 daily values. Regional weather contains 438 of 732 requested cycles.
These counts are from the saved coverage artifact, not a new database query.
Historical publication/revisions remain unverified for several sources; weather
includes hindcast or unverified archive records. A national electricity balance
is not a physical cross-border flow, and an intraday NTC is not a DA capacity.

At an origin $k$, a usable feature is the latest source value with stated
$A_j\le k$, and training labels must also precede the fitting cutoff under the
declared publication policy. The features benchmark explicitly assumes
delivery-day IGI and ENTSO-E day-ahead fundamentals available before DA/IDA1;
its strict alternatives reject unknown clocks. The initial MPC context study
uses 48-hour shifts for French/Swiss realized observations and PSV, and a
nine-day Swiss-border shift with an assumed weekly release. A separate
previous-day experiment is intentionally more permissive. These conventions
are scenario assumptions, not recovered historical publication evidence.

The base feature vector contains local calendar harmonics and eligible price
lags. Fundamentals add NORD load, total generation, solar and wind forecasts.
Context ablations add gas, Milan/Turin weather, French nuclear/hydro and
commercial exchange, Swiss balance, border schedules/NTC/spreads, or reservoir
information. Missing observations remain gaps. Native hourly prices are held
flat when aligned to quarter-hours; interpolation must not invent quarter-hour
target variation. The newer searches separate hourly and quarter-hour labels
before expansion, unlike the explicitly coarse 2024 intraday training extension.

\subsection{Point, spread and distributional prediction problems}
The feature-tree tasks are DA before DA, IDA1--DA spread before DA, that spread
before IDA1, and direct IDA1 before IDA1. Only post-DA origins may use target-day
published DA prices. With $s_t=p_{1,t}-p_{DA,t}$,
\begin{equation}
 \hat p_{1,t}=p_{DA,t}+\hat s_t
 \quad\Longrightarrow\quad
 p_{1,t}-\hat p_{1,t}=s_t-\hat s_t.
\end{equation}
Thus the zero-spread baseline is the published DA curve and is an essential
post-DA comparator. The newer MPC searches also model IDA2, IDA3 and settlement,
use earlier-cleared same-delivery auction curves, and predict spreads relative
to the most recently cleared auction where configured. Adjacent delivery slots
of an already published earlier-auction curve are allowable inputs; later
products of the auction being predicted are not disclosed labels.

For training-standardized predictors $\tilde x_i$, regularized regression solves
\begin{equation}\label{eq:price-regression}
 \min_{\beta_0,\beta}\sum_i\omega_i\rho(y_i-\beta_0-\tilde x_i^{\mathsf T}\beta)
       +\lambda_1\|\beta\|_1+\lambda_2\|\beta\|_2^2.
\end{equation}
Loss/normalization follow each estimator: ridge uses squared loss and $L^2$;
Lasso uses $L^1$ coefficient regularization; elastic net mixes penalties;
Huber estimates robust location. The estimator-specific loss scale matters when
comparing penalty parameters. Spline and polynomial/random-feature variants
change the design matrix. Neural point models, trees and boosted trees use their
specified squared, absolute, Huber or quantile losses. MAE-based selection does
not imply that each model trained on MAE.

The current Huber candidate uses training-only target scaling as well as feature
scaling and inverse-transforms predictions to EUR/MWh. Earlier Huber artifacts
precede that change. Interaction regressions augment noncalendar signals $z$
with $z\sin(2\pi h/24)$ and $z\cos(2\pi h/24)$, allowing time-of-day effects
while retaining a pooled model. Daily-curve inputs concatenate lagged curve
features; delivery-slot fits are a separate representation. Window experiments
restrict eligible training to the last 30, 60 or 90 days versus the full
available history; they change fitting data, not the definition of the target.

In the feature benchmark, Gaussian and Johnson-SU networks minimize negative
log likelihood. With $Z\sim N(0,1)$, the latter has
\begin{equation}
 Y=\xi+\lambda\sinh((Z-\gamma)/\delta),\qquad\lambda,\delta>0,
\end{equation}
with constrained neural parameter maps. Its location $\xi$ is not generally
its mean; actual mean, median, quantiles and PIT are returned separately.
Quantile models minimize pinball loss; mean/median/robust-location outputs are
not silently interchanged. NBEATSx and a bounded local-checkpoint TabPFN pilot
are optional benchmark architectures, not a new probabilistic trajectory law.

\subsection{Residual trajectories, combinations and forecast scores}
The initial MPC ridge workflow sets $\alpha=20$ with a 14-day chronological
calibration tail. A predictor fitted before the tail generates its residuals;
the deployed point predictor is then refitted through the eligible fold cutoff.
For complete local-day residual block $\varepsilon_{d_j,t}$,
\begin{equation}
 \hat\pi^{(j)}_t=\hat\mu_t+\varepsilon_{d_j,t}.
\end{equation}
This preserves empirical within-day error shape. DST matching is explicit;
separate days are sampled independently. It does not learn cross-market,
cross-day or PV--price dependence, and the finite scenario mean need not equal
the point predictor. The basic studies completed 30 chronological fits for
January/August and 27 extended fits including October and coarse 2024 IDA data.

An equal combination averages selected point forecasts. A convex combination
fits earlier out-of-sample predictions $F_{im}$ under
\begin{equation}
 \min_{w\ge0,\ \mathbf1^{\mathsf T}w=1}
 \frac1N\|y-Fw\|_2^2+\lambda\|w\|_2^2,
 \qquad \hat y=\sum_mw_m\hat y_m,
\end{equation}
with scaling/regularization defined by the particular study. Size and penalty
are chosen on an earlier tuning slice; the weights are then refit there before
later evaluation. Local-hour-block combinations select separate pools by time
block. QRA replaces the squared objective by quantile losses. Conformal helpers
use a prior residual quantile for intervals; no universal conditional-coverage
claim follows under drift. Feature-tree QRA/learned combinations are executable
helpers without completed historical evaluation there; the MPC-tree temporal
combination studies have completed saved evaluations.

For price observations and forecasts on a common sample,
\begin{align}
 \mathrm{MAE}&=N^{-1}\sum_i|y_i-\hat y_i|,&
 \mathrm{WAPE}&=100\frac{\sum_i|y_i-\hat y_i|}{\sum_i|y_i|},\\
 \mathrm{rMAE}&=\frac{\sum_i|y_i-\hat y_i|}{\sum_i|y_i-\hat y_i^{\rm naive}|}.
\end{align}
Zero denominators produce undefined relative scores. The feature-tree benchmark
duration-weights within days, then weights days equally; the native-resolution
search reports pooled products. These weightings must not be silently merged.
CRPS evaluates the distribution independently of the point score and cash profit.
Finite quantile-grid integration, analytic Gaussian CRPS and exact empirical
CRPS are labelled separately.

\subsection{Completed feature-tree DA/IDA1 results}
The January--March 2026 validation run
\path{20261023T133005867139Z} completed 90 days, four tasks and 94,996
forecast-statistic records, with no recorded failures. These are not 94,996
independent price observations: several statistics/models share each target.
The study used an 84-day rolling window, saved refit schedule and assumed
availability. Its April--August final-test flag is false in this workflow;
related periods were later inspected elsewhere and are not globally untouched.
\begin{center}\small
\begin{tabular}{llrr}
\toprule Task & Model/features & MAE & RMSE\\\midrule
DA before DA & Previous week & 16.57 & 22.97\\
DA before DA & Elastic net, price/calendar F0 & 11.41 & 16.50\\
DA before DA & Elastic net, fundamentals+gas F2 & 8.11 & 12.57\\
IDA1 spread after DA & Zero spread & 4.90 & 6.61\\
IDA1 spread after DA & LightGBM F2 & 5.02 & 6.66\\
IDA1 spread before DA & LightGBM F2 & 4.98 & 6.65\\
Direct IDA1 after DA & Elastic net F2 & 5.93 & 8.07\\
\bottomrule\end{tabular}\end{center}
All errors are EUR/MWh, daily duration-weighted on the hourly representation.
DA elastic net F2 improves paired MAE by 8.45 EUR/MWh versus previous week;
the saved seven-day block interval is $[-12.06,-5.92]$ for model minus baseline.
This is a combined feature/model contrast, not a gas-only causal effect.
For post-DA IDA1, the zero-spread benchmark is better than these fitted models.

The separate January 5--11 challenger week uses matched dates and daily refits.
DA/post-DA-spread MAEs were ridge 7.27/3.97, LightGBM 6.70/3.62, CatBoost
5.40/3.69, MLP 5.30/3.83, Gaussian network 6.44/3.70, Johnson-SU 5.32/3.95,
NBEATSx 10.24/5.70, TabPFN 5.99/3.99, and equal ridge/MLP 4.91/3.82.
Zero-spread MAE was 3.95. Seven days cannot establish a seasonal architecture
ranking. Separate weekly-refit/full-profile tests gave pooled-versus-slot ridge
DA MAE 12.20/12.87 and spread MAE 4.66/7.70. A matched F0--F4 feature ablation
gave DA MAEs 10.74, 8.51, 12.20, 11.29, 11.29: more context did not uniformly
help. These different refit/design settings prevent cross-run architecture
attribution. The October 26 compact-design four-day pilot is also complete,
but is a workflow check rather than a replacement for full-quarter validation.

\subsection{MPC-tree price studies and selection-fixed follow-up}
The first chronological ridge study reports the following representative
point/distribution comparison on common days:
\begin{center}\small
\begin{tabular}{llrrr}
\toprule Period/market & Model & Days & MAE & CRPS\\\midrule
January DA & Price-only ridge & 31 & 10.260 & 9.183\\
January DA & Fundamentals ridge & 31 & 7.545 & 6.915\\
August DA & Price-only ridge & 31 & 17.123 & 14.542\\
August DA & Fundamentals ridge & 31 & 19.011 & 12.192\\
January settlement & Fundamentals ridge & 31 & 47.670 & 35.065\\
August settlement & Fundamentals ridge & 31 & 73.339 & 45.169\\
\bottomrule\end{tabular}\end{center}
Fundamentals improve January point accuracy but worsen August DA point MAE
while improving its CRPS. Distribution and point metrics therefore support
different conclusions. In the separate European-context January DA ablation
on 30 common days, price-only/fundamentals/combined MAEs are 10.427/7.715/7.285;
coverage and training samples differ from the 31-day study, so those numbers
are not a single controlled sequence of gains.

Expanded searches examine native resolution, daily curves, slot models,
earlier-cleared auctions, spread targets, French/Swiss context, Terna load,
recent windows and robust regression. All 40 captured run manifests, their
status and saved record counts are listed in Appendix~\ref{app:price-inventory};
the accompanying CSV retains every saved candidate score, not just winners.
The central later-evaluation evidence is:
\begin{center}\small
\begin{tabular}{llrrrr}
\toprule Market & Earlier-selected model & Later MAE & WAPE & rMAE & Products\\\midrule
IDA1 & Spread/combined LightGBM L1 & 6.88 & 6.10\% & 0.498 & 8160\\
IDA2 & Cleared-only LightGBM L1 & 7.12 & 6.26\% & 0.472 & 8064\\
IDA3 & Cleared/combined Lasso 1 & 10.54 & 9.08\% & 0.534 & 4176\\
\bottomrule\end{tabular}\end{center}
These are the \emph{first-ranked} candidates from June--August 2025 selection,
followed on shared complete days from September 3--November 30 with 13 weekly
refits and a one-day gap. Products and label publication must precede each fit.
The full follow-up contains 27 fits, including other earlier-selected candidates
and ridge references. Choosing its best row after inspection would be a new
selection. Historical availability remains assumed and the broader research
history prevents a globally untouched-test claim.

For native quarter-hour DA, previous-day IDA1/IDA2 prices at 24/48-hour lags
were tested across 35 candidates and matched feature sets. The February--April
validation minima were 9.46\% WAPE with IDA inputs versus 9.64\% without, on
82 common days. The earlier top candidate, fundamentals interaction ridge 200,
was then followed on 107 complete days in May 3--August 31, with 18 weekly
refits and 10,272 products: MAE changes from 16.314 to 16.021 EUR/MWh and
WAPE from 11.23 to 11.03\%. All ten matched follow-up fits completed.
The small paired benefit is distinct from selecting a different model after
seeing later scores.

Temporal combinations rank candidates on February 1--March 15, tune on
March 16--31 and evaluate April 1--30, 2026. Equal top-ten ensembles give
IDA1/IDA2/IDA3 MAEs 5.989/7.836/13.525 EUR/MWh and WAPEs
4.94/6.44/12.10\%. Convex stacking gives 6.040/7.838/13.563; local-hour
combinations give 5.995/7.809/13.525. More complex combination did not
uniformly improve this previously inspected later slice.

Matched French-balance and Terna-load ablations show heterogeneous, often tiny
median MAE changes across 27 candidates per task/feature pair. Swiss-reservoir
median changes range from $-0.004$ to $-0.017$ EUR/MWh in the saved comparisons.
Those exploratory multiple comparisons do not establish an economically
meaningful causal benefit. Pairing training rows as well as test products is
necessary to separate added information from coverage changes.

The 96-fit recent-window DA search is complete: its lowest validation WAPE is
9.01\% for 90-day fundamentals interaction elastic net, versus 9.46\% for the
previous best full-history/IDA configuration. These are selected minima on the
same period. The paired later rolling-window comparison is still recorded as
running at capture, so no later advantage is asserted. The bidirectional
earlier-auction-neighbor search is also incomplete; its offsets are $-60$,
$-15$, $15$ and $60$ minutes, with saved boundary handling.

The first DA Huber search completed 56 fits but logged iteration-limit warnings;
its scores are provisional. A target-scaled 1,000-iteration probe still warned.
The 144-fit IDA Huber search has no warning in its relevant old shared-log part,
but lacks candidate-level convergence evidence. Its minimum Huber WAPEs were
4.89/5.48/9.18\%, versus Lasso minima 4.74/5.49/9.18\%. Current code records
warnings per candidate and excludes warned fits from ranking. Completed output
is not proof of numerical convergence, and current warning handling cannot be
retroactively assumed for old results.

\clearpage
\section{Real-time research extensions: mathematical specification}
\label{sec:research-extensions}

\subsection{Forecast-economic control and pathwise monetary risk}
The MPC sibling extends the common-control model with settlement forecasts,
adaptive PV-history recourse, dependence transformations, scenario reduction,
continuation values and controlled information/asset experiments. Its launcher
uses continuous replay: each case carries its own actual SOC, commitments,
degradation and fee ledger. Fixed-ledger replay remains an older diagnostic.

Let $n(s,t)$ own controls and let $a_{n,t}\in[0,1]$ be PV utilization, avoiding
confusion with the earlier curtailment-power variable $u$. Define
\begin{align}
 g_{s,t}&=a_{n(s,t),t}V_{s,t}+d_{n(s,t),t}-c_{n(s,t),t},\\
 b^+_{s,t}-b^-_{s,t}&=\Delta_tg_{s,t}-Q_t^{\rm sub},\qquad b^\pm_{s,t}\ge0,\\
 R_s&=\sum_t(\hat\lambda^+_{s,t}b^+_{s,t}-\hat\lambda^-_{s,t}b^-_{s,t}),
 \qquad L_s=D_s+I_s-R_s.
\end{align}
The split variables are bounded by
$b^+_{s,t}\le\max(0,P^{\rm exp}\Delta_t-Q_t^{\rm sub})$ and
$b^-_{s,t}\le\max(0,P^{\rm imp}\Delta_t+Q_t^{\rm sub})$.
If directional prices differ, a sign binary enforces exclusivity. With equal
prices simultaneous split slack cancels from settlement; a positive volume
penalty removes it. Reported imbalance is computed from physical exchange.
Here $D_s$ is executed-path planned wear, $I_s$ import cost and $K_s$ terminal
cost (possibly negative continuation value). All battery transitions, initial
state, POI, power/SOC and terminal constraints apply on each path. There are
\emph{no new auction orders} in this physical-control problem. Forecast economics
minimizes
\begin{align}\label{eq:research-mpc}
 \min\quad &\sum_sw_s\left[D_s+I_s+K_s-R_s+\kappa_{\rm econ}
                 \sum_t(b^+_{s,t}+b^-_{s,t})\right]
 +\lambda_{\rm risk}\left[\eta+
        \frac1{1-\alpha}\sum_sw_s\xi_s\right],\\
 &\xi_s\ge L_s-\eta,\quad\xi_s\ge0.\notag
\end{align}
$\eta$ is unrestricted, because monetary loss can be negative. This differs
from nonnegative absolute-volume CVaR in the auction model. Risk excludes the
artificial imbalance regularizer, terminal salvage and uncertain cash on
already submitted unpriced auctions. Research economic cases default to
$\kappa_{\rm econ}=0$, independently of the 500 EUR/MWh price-blind penalty;
the risk-neutral baseline has $\lambda_{\rm risk}=0$.

For unequal directional prices, sign exclusivity prevents artificial positive
and negative settlement at once. Equal prices reduce cash to
$\hat\lambda_{s,t}(b^+_{s,t}-b^-_{s,t})$; slack in the split cannot create cash.
Without a positive absolute-volume coefficient, report reconstructed
$|\Delta_tg_{s,t}-Q_t^{\rm sub}|$, not an arbitrary slack split. Nonnegative
import costs use an epigraph tight at an optimum with positive coefficient.
Only the first control executes; physical protection and rejected-plan fallbacks
are logged. Artificial risk, tracking and terminal terms never enter cash profit.

\subsection{Observable-history recourse and the value of dependence}
Ordinary deterministic/ensemble MPC shares all horizon controls and energy
across paths. Adaptive MPC shares the first command but allows future controls
to differ only after distinguishable observed PV histories. With lag $\ell\ge1$
and quantization width $q$, a branch key is
\begin{equation}
 H_{s,t}=\left(n(s,t-1),
       \left\lfloor V_{s,0:t-\ell}/q\right\rfloor\right),
 \qquad H_{s,t}=H_{r,t}\Rightarrow n(s,t)=n(r,t).
\end{equation}
Zero quantization uses exact observations; configured nonbranch times retain
the previous owner. The previous node prevents recombining different battery
histories. Future prices and member labels do not enter the observation key.
The prepared longer study uses hourly branching (four quarters), one-quarter
observation lag and 0.1 MW bins. Exact small-ensemble prefixes can identify a
member early, so finite-ensemble recourse is still an approximation to an
operational information process.

Raw member trajectories preserve their supplied serial structure. Independent
per-lead shuffling permutes the \emph{joint} PV/price sample at each lead:
\begin{equation}
 (V'_{s,t},\lambda'_{s,t})=(V_{\pi_t(s),t},\lambda_{\pi_t(s),t}).
\end{equation}
Equal weights and within-lead joint samples are retained while serial dependence
is disrupted. Common open-loop controls with an additive risk-neutral objective
are invariant to this operation: their feasible POI support and each lead's
expected objective are unchanged. Adaptive recourse and pathwise monetary risk
provide mechanisms by which dependence can matter; this invariance is a useful
negative control, not a claim that dependence never matters for storage.

ECC uses calibrated marginal samples $z_{(1),t}\le\cdots\le z_{(M),t}$ and the
raw member's rank $r_{s,t}$:
\begin{equation}
 V^{\rm ECC}_{s,t}=z_{(r_{s,t}),t}.
\end{equation}
Seeded tie handling is explicit. Price row identities are retained, and this
does not estimate a new joint PV--price copula. The construction follows the
rank-rearrangement principle of ECC~\cite{ecc}; its marginal calibration must
be supplied separately. Calibrated ECC and its shuffled counterpart have the
same marginal scores, allowing a dependence contrast under matched control.

The reviewed-target affine bias/spread calibration fits January 1--14 and freezes
parameters for January 15--16 and August 2--3. Its two-parameter form in MW is
\begin{equation}
 V'_{s,t}=\clip\{\bar V_t+aP_{\rm AC}+c(V_{s,t}-\bar V_t),0,P_{\rm AC}\},
\end{equation}
where normalized bias $a\in[-0.3,0.3]$ and spread multiplier
$c\in[0.25,3]$ minimize training mean CRPS. The ensemble-mean coefficient
is fixed to one; monthly night masks are applied after clipping.
On the August two-day transfer, raw versus calibrated CRPS in MW was
0.02554/0.06944 (farm 1), 0.14332/0.15116 (345), 0.19844/0.22464 (901), and
0.17979/0.29941 (910). Calibration worsened all four marginal scores there.
Farm 910's 192 execution targets comprise 132 historical-GHI reconstructions
and 60 policy night zeros, so that transfer is not independent meter validation.

\subsection{Decision-aware scenario reduction}
The reduction preserves original complete paths and redistributes probability
mass; it does not take a pointwise minimum. Let $\phi_s$ concatenate scaled
features of PV, ramps, cumulative energy deficit/balance, capacity/POI
exceedance proxies, price and uncontrolled settlement exposure. Each block
is divided by its unit scale and multiplied by the square root of its weight
over block length. The implementation greedily approximates
\begin{equation}
 \min_{\mathcal M:\,|\mathcal M|=K}\sum_sw_s\min_{j\in\mathcal M}
       \|\phi_s-\phi_j\|_2,
 \qquad \tilde w_j=\sum_{s:\,j=\operatorname{assign}(s)}w_s.
\end{equation}
Representatives are source members and retain positive mass, including tie
handling for duplicate paths. Features use only current forecasts, known
commitments and asset parameters. The distance is a heuristic relevance metric,
not a bound on lost MILP value, tail risk or physical feasibility on discarded
scenarios. This approximate reduction must not be confused with the exact
PV-envelope theorem for the special auction formulation.

\subsection{Conditional terminal continuation value}
A finite daily horizon can empty or retain storage for reasons unrelated to
future opportunities. The research continuation model uses available nominal
PV/settlement forecasts and \emph{known} commitments after the MPC horizon.
On an energy grid $\mathcal E$, define
\begin{equation}
 W_t(e)=\max_{e'\in\mathcal E\ {
m feasible}}
       \{r_t(e,e')+W_{t+1}(e')\},\qquad W_{T'}(e)=v_F e.
\end{equation}
For a transition $h=e'-\rho_te$,
$c=\pos{h}/(\eta_c\Delta_t)$ and
$d=\eta_d\pos{-h}/\Delta_t$, subject to power limits.
Given nominal PV, feasible grid power ranges from
$\max(d-c,-P^{\rm imp})$ to $\min(d-c+V_t,P^{\rm exp})$.
The one-step reward maximizes signed settlement minus import cost and
throughput wear; a piecewise-linear maximum can be checked at range endpoints,
zero grid flow and the commitment kink. Idle with interpolated continuation
is also considered when decay leaves the discrete grid.

The resulting curve is linearly interpolated and enters MPC as
$K_s=-W_T(e_{s,T})$ using an exact PWL graph for that supplied curve.
The curve itself is approximate: nominal uncertainty, discretized energy,
throughput wear without Xu memory, and no unmodeled future auction recourse.
It is neither a learned lifetime value nor realized income. A continuation
gain must be compared on matched terminal settings and final inventory.

\subsection{Information/access and physical-flexibility experiments}
The information/access factorial crosses latest versus frozen-initial
information $I\in\{0,1\}$ with DA-only versus all-auction access
$A\in\{0,1\}$. For matched realized profits $J_{IA}$, the interaction is
\begin{equation}
 \Gamma=J_{11}-J_{10}-J_{01}+J_{00}.
\end{equation}
A positive value indicates complementarity on those cases and dates; it is
not identified if cells differ in assets, targets or forecast families.
Paired reports require all four cells and preserve their independent state.
Separate contrasts double energy capacity, halve power or halve POI capacity;
crossing each with marginal calibration tests whether the calibration effect
depends on flexibility. For asset setting $b$, the difference-in-differences is
$(J_{{\rm cal},b}-J_{{\rm raw},b})-(J_{{\rm cal},0}-J_{{\rm raw},0})$.
These are configured asset counterfactuals, not measured farm upgrades.

\subsection{Oracle certification, time budgets and numerical recovery}
The main-tree oracle described earlier shares a ten-second budget and requires
its stage-one certificate. The MPC sibling additionally supports separate
ten-second stage budgets and a declared 1\% time-limit fallback. Let $U$ be
independently reconstructed feasible absolute imbalance and $L$ the lower bound.
The fallback requires $(U-L)/U\le0.01$ for positive $U$, with explicit near-zero
handling; the default target solver gap remains $10^{-4}$. Refinement is bounded by
\begin{equation}
 A\le U(1+r)+b.
\end{equation}
With an approximate incumbent this is an incumbent cap, not a certified
$A_{\min}$ cap. Possible excess over the true minimum is at most
$U-L+rU+b$, subject to the recorded numerical tolerances. Invalid refinement
retains the validated stage-one action; no valid action triggers tracking
fallback. Those are different events and are counted separately.

Optional numerical retry tightens feasibility tolerance and sets numerical focus
within the remaining economic-stage time budget. It cannot change physical
constraints or weaken the imbalance-cap check. The completed 104-job matrix
predates these two added retry cases: current configuration has 28 cases,
whereas the completed matrix has 26. The current two-window/four-farm plan
therefore describes 224 jobs, not 224 completed results. Timing reports separate
model construction, native solver runtime, validation and end-to-end replay.

\clearpage
\section{Performed trading experiments and numerical evidence}
\label{sec:empirical-results}

\subsection{Independent-day participation and continuous monthly auction studies}
The completed January 2025 farm-910 noon-boundary participation run contains
15 integrated profiles on 31 independently initialized days and a sequential
all-auction reference. Its all-auction integrated profile ranks first with
mean EUR 3,520.06/day, total EUR 109,121.72, daily population standard deviation
EUR 1,187.52, mean absolute imbalance 1.69984 MWh/day and mean ending SOC
28.25\%. These are price-oracle operating results, not forecast-price earnings
or battery-only value. The next-ranked IDA1+IDA2+IDA3 profile averages
EUR 2,893.38/day. Independent daily initialization and fee-state reset prevent
interpreting their sum as one continuous asset trajectory.

The later continuous participation study instead carries state separately for
each of 15 profiles plus sequential and perfect-PV references. Four farms and
two forecast variants have saved full-January comparisons. Every strategy has
its own SOC, degradation memory, commitments and fee state; the best profile
is an ex-post ranking on January. Appendix~\ref{app:participation} gives all
68 matched profile rows. For farm 1, all-auction deterministic/ensemble profit
is EUR 10,216.33/10,197.54, versus sequential EUR 1,848.22/1,798.72 and
perfect-PV EUR 10,582.72. The large integrated/sequential gap includes
perfect-price cross-auction spread opportunities; it is not purely the value
of additional PV forecast information.

The separate mixed-model full-month comparisons below use January 1--31 and
August 2--31, 2026. Each month starts independently at its configured state;
August does not carry January SOC or simulate February--July. The absent
July 31 ensemble cycle excludes August 1. Farm 345 uses the declared pvlib
override and the other farms retained TCNs. These earlier target/model policies
differ from the new reviewed-GHI registry and from the continuous-participation
run. The table consistently uses the named \texttt{mixed\_models} comparison
family rather than mixing its revisions with older all-TCN summaries.
\begingroup\small
\begin{longtable}{llrrrrr}
\toprule
Farm/month & Forecast & Profit EUR & Abs. dev. MWh & Wear EUR & Curt. MWh & End SOC \%\\
\midrule\endhead
1/01 & Det. & 10,319.57 & 0.665 & 452.26 & 0.267 & 44.14\\
1/01 & 51 paths & 10,054.81 & 0.226 & 531.35 & 1.835 & 50.17\\
1/01 & Perfect PV & 10,582.72 & 0.000 & 224.23 & 0.000 & 50.00\\
1/08 & Det. & 28,373.64 & 3.621 & 900.23 & 2.219 & 48.27\\
1/08 & 51 paths & 26,986.34 & 0.816 & 883.62 & 11.013 & 49.32\\
1/08 & Perfect PV & 28,991.90 & 0.000 & 768.58 & 0.000 & 50.00\\
345/01 & Det. & 20,134.68 & 48.063 & 689.45 & 0.006 & 19.36\\
345/01 & 51 paths & 19,643.28 & 18.672 & 922.97 & 2.959 & 16.67\\
345/01 & Perfect PV & 20,956.88 & 0.000 & 283.99 & 0.000 & 50.00\\
345/08 & Det. & 50,554.56 & 45.594 & 922.19 & 0.310 & 7.58\\
345/08 & 51 paths & 49,053.88 & 10.347 & 1,428.87 & 14.229 & 7.07\\
345/08 & Perfect PV & 52,615.69 & 0.000 & 967.83 & 0.000 & 50.00\\
901/01 & Det. & 53,451.80 & 14.022 & 1,920.99 & 3.963 & 93.00\\
901/01 & 51 paths & 49,170.14 & 2.071 & 1,831.25 & 42.087 & 93.00\\
901/01 & Perfect PV & 56,170.13 & 0.000 & 533.91 & 0.000 & 50.00\\
901/08 & Det. & 177,589.29 & 38.691 & 2,646.46 & 7.991 & 52.14\\
901/08 & 51 paths & 159,278.29 & 3.251 & 2,627.54 & 128.051 & 53.31\\
901/08 & Perfect PV & 181,957.15 & 0.000 & 1,935.66 & 0.000 & 50.00\\
910/01 & Det. & 77,704.74 & 27.099 & 3,957.66 & 14.125 & 48.68\\
910/01 & 51 paths & 74,492.63 & 10.151 & 3,983.35 & 45.726 & 49.32\\
910/01 & Perfect PV & 83,258.82 & 0.000 & 1,295.10 & 0.000 & 50.00\\
910/08 & Det. & 222,503.76 & 422.351 & 3,454.07 & 0.414 & 12.35\\
910/08 & 51 paths & 216,492.01 & 189.140 & 5,407.93 & 52.514 & 7.66\\
910/08 & Perfect PV & 239,605.95 & 0.000 & 4,839.14 & 0.000 & 50.00\\
\bottomrule
\end{longtable}
\endgroup

In these saved cases the exact 51-path policy reduces absolute imbalance but
also reduces cash profit relative to deterministic planning and often increases
curtailment. This agrees with the common-control exact-delivery envelope's
conservative feasibility mechanism; it does not establish that probabilistic
forecasting intrinsically reduces value. The constraints, execution rule,
calibration and objective determine how uncertainty is used. Farm 910's large
August imbalance belongs to the older target/model setting, not to the later
fully reconstructed two-day MPC target. Final SOC differences are shown because
cash profit contains no terminal salvage valuation.

The envelope equivalence experiment compares 37 matched auction states across
August 2--11, farm 1, then two independent ten-day replays. Both saved replays
have profit EUR 10,366.671392, absolute imbalance 0.0123333 MWh, wear
EUR 283.344853 and final SOC 74.086594\%, agreeing to numerical precision.
The recovered-curtailment diagnostic differs because it is relative to the
retained plan representation. This does not contradict equal actual flows.
The fixed-state equivalence follows the stated theorem; equality of independent
replay schedules is empirical and can change under alternative optima.

The 100/500 EUR/MWh stochastic-imbalance sensitivities, ensemble-mean versus
full paths, selected 10/20-day participation contrasts and PV-only battery
counterfactuals have separate comparison artifacts. They are distinct
objectives/samples and are not pooled with the exact-delivery monthly table.
Their numerical records are preserved in the sensitivity appendix; no artificial
planning penalty is relabelled as realized settlement.

\subsection{Original current-PV-known two-day MPC diagnostic}
The farm-1 January 15--16 TCN study disclosed current-interval realized PV to
MPC but kept future PV forecast-only. On a frozen common ledger, tracking
profit was EUR 538.25 with 0.217316 MWh absolute imbalance. Deterministic MPC
earned EUR 550.53 with 0.236596 MWh; ensemble MPC EUR 554.73 with 0.258667 MWh.
Their price-oracle variants earned EUR 541.99 and EUR 541.41. All MPC cases
ended at 50\% SOC versus tracking 50.677\%. In the continuous deterministic
case, profit was EUR 556.16 and imbalance 0.361619 MWh; the EUR 17.91 gain
included EUR 24.30 more auction revenue from later changed submissions.
The 960 recorded MPC decisions had no tracking fallback or timeout. These
figures are a software/development diagnostic with retrospective information,
not the forecast-current-PV reviewed-target matrix below.

\subsection{Reviewed-target October week}
For October 1--7, 2025, deterministic penalty MPC versus its own continuous
tracking baseline increased cash profit by EUR 44.53, 675.75, 1,157.14 and
1,408.45 for farms 1, 345, 901 and 910. Portfolio absolute imbalance increased
89.91 MWh. The saved independent audit reconstructs 2,688 MPC decisions with
no physical/accounting violations, tracking fallback or timeout. It uses a
$10^{-6}$ reconstruction tolerance; this is numerical consistency evidence,
not forecast or economic validation. Ending battery energy differs, and MPC
has a soft inventory objective absent from tracking. The result therefore
does not isolate control intelligence from inventory preference and information.

\subsection{Completed reviewed-target January matrix}
The January 15--16, 2026 matrix completed 104/104 jobs: 26 cases by four farms,
with a separate tracking baseline for each case. Each target has 192 MPC
decisions, hence 19,968 target MPC decisions in total. Saved resolved settings
use 50\% initial SOC, a 50\% soft terminal target at 50 EUR/MWh, hourly
recourse branching, one-quarter observation lag and 0.1 MW PV bins. The
reduced case retains eight paths with seed 17; the monetary-risk case uses
$\alpha=0.9$, $\lambda_{\rm risk}=0.5$. Continuation replaces the soft target
with a six-hour, 51-energy-point nominal value curve and zero final salvage.
Auction solves have 120 seconds and MPC stages 10 seconds each. The frozen
\path{january_case_settings.json} retains all 26 farm-1 resolved configurations
and their source hashes; farm-specific asset parameters remain in the run plan.
The following values
sum both days and all four farms; they are \emph{not} mean daily gains. Positive
imbalance/wear means more absolute deviation/more wear than that case's baseline.
\begingroup\small
\begin{longtable}{lrrrr}
\toprule
Case & Profit EUR & Abs. dev. MWh & Wear EUR & End energy MWh\\
\midrule\endhead
adaptive\_continuation & 331.45 & 27.795 & 297.96 & -0.122\\
adaptive\_ecc & 302.05 & 24.400 & 256.79 & -0.122\\
adaptive\_monetary\_risk & 94.83 & 21.076 & 284.63 & -0.122\\
adaptive\_raw & 301.73 & 23.550 & 256.51 & -0.122\\
adaptive\_reduced & -130.31 & 32.188 & 564.56 & -0.122\\
adaptive\_shuffled & 332.26 & 26.487 & 345.75 & -0.122\\
calibrated\_double\_energy & 497.05 & 30.093 & 289.66 & -0.122\\
calibrated\_half\_poi & 321.95 & 24.657 & 255.43 & -0.122\\
calibrated\_half\_power & 290.14 & 15.504 & 118.23 & -0.122\\
calibrated\_shuffled & 355.97 & 26.983 & 342.83 & -0.122\\
causal\_auctions & 300.53 & 21.665 & 303.46 & -0.122\\
causal\_fundamentals & 249.38 & 27.548 & 542.42 & -0.122\\
da\_only\_frozen & 282.72 & 20.025 & 110.42 & -0.505\\
da\_only\_latest & 260.57 & 22.941 & 220.69 & -0.505\\
deterministic\_economic & 98.10 & 12.294 & -103.12 & -0.122\\
deterministic\_oracle\_gap\_fallback & 218.42 & 5.573 & 16.64 & 0.363\\
deterministic\_penalty & 123.66 & 5.707 & 17.82 & 0.337\\
deterministic\_price\_oracle & 218.42 & 5.573 & 16.64 & 0.363\\
double\_battery\_energy & 478.19 & 29.771 & 296.91 & -0.122\\
ensemble\_economic & 177.81 & 20.784 & 187.46 & -0.122\\
ensemble\_oracle\_gap\_fallback & 8.58 & 10.565 & 307.30 & -0.122\\
ensemble\_penalty & 74.81 & 9.114 & 186.09 & -0.122\\
ensemble\_price\_oracle & 8.58 & 10.565 & 307.30 & -0.122\\
frozen\_information & 233.48 & 21.544 & 155.66 & -0.098\\
half\_battery\_power & 294.71 & 15.385 & 121.86 & -0.122\\
half\_poi\_capacity & 309.85 & 23.568 & 245.70 & -0.122\\
\bottomrule
\end{longtable}
\endgroup

The deterministic penalty, ensemble penalty and adaptive-raw gains correspond
to EUR 61.83, 37.40 and 150.86 per portfolio day. Their absolute-imbalance
increases are 2.85, 4.56 and 11.78 MWh per day. Adaptive reduction gives a
negative aggregate cash difference of EUR 130.31, whereas raw adaptive gives
a positive EUR 301.73. That is evidence that this particular heuristic/sample
can alter economic decisions substantially, not a general reduction theorem.
Calibrated ECC and raw adaptive have nearly equal total cash differences here,
while shuffling and continuation can change them. The whole two-day sample
is too small to rank these policies reliably or estimate tail risk.

Some rows change objectives, asset ratings, forecast information or auction
access, so subtraction between arbitrary rows does not isolate a single
mechanism. For a pure recourse contrast use matched common/adaptive economic
objectives and scenarios; for dependence use matched marginals; for flexibility
use its declared factorial contrasts. Each baseline belongs to its own case.
Per-farm differences, rather than only portfolio totals, appear in
Appendix~\ref{app:mpc-farms} and the evidence CSV retains absolute target and
baseline metrics, auction cash, settlement, charge/discharge and curtailment.

Using absolute target profits (not differences from changing tracking baselines),
the four-cell information-by-auction-access interaction is EUR 206.84 summed
over the portfolio and two days. The calibration-by-flexibility interactions
are EUR 18.54 for doubled battery energy, EUR $-4.89$ for halved power and
EUR 11.77 for halved POI capacity. These are descriptive contrasts with common
targets and dates; two calendar days cannot support a reliable confidence
interval or a claim about seasonal interactions.

The saved all-job audit reports zero physical/accounting/artifact violations.
No target used a tracking fallback or solver timeout. Eighteen economic oracle
refinements across the ordinary and fallback cases were rejected by validation;
validated stage-one actions were retained. These are not 18 invalid accepted
actions. Accepted cap residuals stayed within the documented allowance. The
older raw-target 32-job batch had 18 excessive cap residuals in two oracle jobs
and remains superseded evidence; corrected eight-job oracle results have their
own audit. Numerical correctness is stated for each audited batch, not inferred
for every historical output from current software tests.

\subsection{Reviewed-target August farm-910 comparison}
On August 2--3, ensemble penalty MPC has cash gain EUR 96.812, extra absolute
imbalance 7.542783 MWh and wear reduction EUR 163.559 versus tracking.
Adaptive calibrated ECC gains EUR 271.803, with extra imbalance 26.960265 MWh,
extra wear EUR 165.250 and four physical safety overrides. Ensemble price oracle
loses EUR 62.935, with extra imbalance 7.501351 MWh, wear reduction EUR 119.006
and 40 time-limited decisions. The 1\% fallback variant has essentially the
same realized outcome. Each has 192 MPC decisions and no tracking fallback;
the saved audit reports no physical/accounting violations. All 192 target
intervals are reconstructed daylight or policy night zeros. This is a
counterfactual control comparison, not measured plant-performance evidence.

\subsection{What remains unfinished}
The main-tree May all-farm study is explicitly interrupted with no completed
comparison mode. Its successful preflight is not a controller result. The
prepared annual deterministic study and full January/August ensemble studies
still require complete data and final model integration. Reviewed PV inputs
end in August; the missing July 31 ensemble cycle remains material for August 1.
The current 28-case research plan adds numerical-retry cases beyond the
completed 26-case matrix. Four-cycle neural PV training, recent-window rolling
price evaluation and bidirectional-neighbor price search are partial at capture.
Mayer-style marginal post-processing has implemented formulations but no
completed numerical output in the inspected trees. A future run must not use
this document's snapshot label as evidence of present process status.

\section{Interpretation and current scientific focus}
The completed studies support a working chronological decision framework and
several useful empirical observations: price-history/fundamental models can
beat persistence on specific samples; calibrated marginals and online weights
do not improve every season or gate; strict ensemble delivery can lower
imbalance while sacrificing cash profit; and physical MPC can change both
settlement and subsequent auction decisions. They do not yet establish a
deployable annual strategy, an investment return, or a universal architecture
ranking.

The current work focuses on forecast-vintage-aligned neural PV models,
selection-fixed native-resolution price comparisons, coherent scenario
dependence under adaptive control, and the interaction of information with
auction access and storage/POI flexibility. Terminal continuation, monetary
risk and reduction remain controlled extensions. A calibrated marginal CDF,
an exact finite-scenario solve and a numerically reconciled replay answer
different questions; none alone establishes real-world reliability.

Statistical units are calendar blocks, shared across farms where prices/weather
are common. The four farms are not independent replications of the market.
Two-day confidence intervals or extreme-tail estimates are not informative;
block bootstrap intervals on larger samples remain conditional on selection,
coverage and target construction. Observed/synthetic target performance, native
product regimes, forecast issue availability, optimization gaps, fallback rates
and final inventory must accompany any future claim of improvement.

\appendix
\section{Complete continuous participation results}
\label{app:participation}
These are four independent January 2026 comparisons. D/E denote updating
deterministic/frozen ensemble plus updating deterministic forecasts. Profit is
EUR, absolute imbalance is MWh and ending SOC is percent. All rows cover
January 1--31; each profile carries its own state. Different studies elsewhere
in this report must not be merged by farm/date alone.
\subsection{Farm 1}
\begingroup\small
\begin{longtable}{lrrrrrr}
\toprule
Profile & Profit D & Profit E & Abs. D & Abs. E & SOC D & SOC E\\
\midrule\endhead
DA & 1,648.99 & 1,748.59 & 9.651 & 2.301 & 7.00 & 36.56\\
IDA1 & 1,665.96 & 1,730.08 & 10.026 & 1.852 & 34.19 & 38.71\\
IDA2 & 1,779.95 & 1,700.89 & 7.800 & 1.638 & 50.00 & 52.06\\
IDA3 & 1,643.95 & 1,553.62 & 5.190 & 1.137 & 33.72 & 43.43\\
DA+IDA1 & 4,974.58 & 5,017.71 & 9.186 & 1.656 & 7.00 & 38.88\\
DA+IDA2 & 6,193.08 & 6,135.73 & 7.064 & 1.505 & 38.56 & 52.28\\
DA+IDA3 & 5,101.96 & 5,170.96 & 6.531 & 2.130 & 33.57 & 43.76\\
IDA1+IDA2 & 5,613.30 & 5,576.99 & 7.469 & 1.532 & 38.66 & 52.50\\
IDA1+IDA3 & 5,018.40 & 5,136.20 & 6.317 & 1.466 & 41.68 & 43.65\\
IDA2+IDA3 & 4,829.37 & 4,772.28 & 5.762 & 1.562 & 33.72 & 43.41\\
DA+IDA1+IDA2 & 7,644.73 & 7,596.64 & 7.636 & 1.506 & 39.64 & 52.28\\
DA+IDA1+IDA3 & 7,585.70 & 7,685.52 & 6.321 & 1.465 & 44.35 & 43.60\\
DA+IDA2+IDA3 & 8,455.24 & 8,446.08 & 5.817 & 1.565 & 33.43 & 43.49\\
IDA1+IDA2+IDA3 & 7,967.28 & 7,963.91 & 5.805 & 1.534 & 33.50 & 43.49\\
DA+IDA1+IDA2+IDA3 & 10,216.33 & 10,197.54 & 5.818 & 1.577 & 33.50 & 43.49\\
Sequential & 1,848.22 & 1,798.72 & 6.430 & 1.896 & 33.38 & 43.49\\
Perfect PV & 10,582.72 & 10,582.72 & 0.000 & 0.000 & 50.00 & 50.00\\
\bottomrule
\end{longtable}
\endgroup

\subsection{Farm 345}
\begingroup\small
\begin{longtable}{lrrrrrr}
\toprule
Profile & Profit D & Profit E & Abs. D & Abs. E & SOC D & SOC E\\
\midrule\endhead
DA & 6,494.09 & 4,796.74 & 3.812 & 1.148 & 59.96 & 62.23\\
IDA1 & 6,624.83 & 4,793.22 & 4.203 & 0.882 & 73.72 & 92.04\\
IDA2 & 6,773.60 & 4,858.31 & 3.427 & 0.735 & 67.79 & 91.81\\
IDA3 & 5,490.12 & 3,957.99 & 1.136 & 0.095 & 46.57 & 49.05\\
DA+IDA1 & 11,835.88 & 10,095.03 & 3.817 & 0.732 & 55.20 & 75.39\\
DA+IDA2 & 13,710.46 & 11,920.12 & 3.323 & 0.660 & 67.83 & 92.04\\
DA+IDA3 & 12,101.25 & 10,462.54 & 2.985 & 1.029 & 75.92 & 89.22\\
IDA1+IDA2 & 12,844.82 & 11,056.18 & 3.789 & 0.685 & 67.86 & 92.04\\
IDA1+IDA3 & 12,140.08 & 10,382.20 & 2.549 & 0.845 & 75.72 & 89.26\\
IDA2+IDA3 & 11,691.12 & 9,894.06 & 2.405 & 0.836 & 76.00 & 89.22\\
DA+IDA1+IDA2 & 16,126.41 & 14,285.36 & 3.661 & 0.689 & 67.82 & 92.04\\
DA+IDA1+IDA3 & 16,201.23 & 14,471.14 & 2.622 & 0.891 & 75.72 & 89.26\\
DA+IDA2+IDA3 & 17,477.33 & 15,727.12 & 2.448 & 0.832 & 75.99 & 89.23\\
IDA1+IDA2+IDA3 & 16,701.51 & 14,941.97 & 2.553 & 0.830 & 76.00 & 89.26\\
DA+IDA1+IDA2+IDA3 & 20,306.16 & 18,555.04 & 2.447 & 0.780 & 76.00 & 89.22\\
Sequential & 6,916.97 & 5,180.87 & 2.676 & 0.837 & 75.91 & 89.22\\
Perfect PV & 20,956.88 & 20,956.88 & 0.000 & 0.000 & 50.00 & 50.00\\
\bottomrule
\end{longtable}
\endgroup

\subsection{Farm 901}
\begingroup\small
\begin{longtable}{lrrrrrr}
\toprule
Profile & Profit D & Profit E & Abs. D & Abs. E & SOC D & SOC E\\
\midrule\endhead
DA & 16,467.77 & 12,232.49 & 48.292 & 3.262 & 7.00 & 47.53\\
IDA1 & 16,963.86 & 11,909.29 & 47.777 & 4.615 & 30.72 & 93.00\\
IDA2 & 17,212.29 & 12,156.39 & 33.781 & 5.321 & 61.64 & 93.00\\
IDA3 & 12,874.63 & 9,399.95 & 16.635 & 1.258 & 36.79 & 76.04\\
DA+IDA1 & 31,169.82 & 26,579.68 & 45.542 & 3.173 & 7.00 & 93.00\\
DA+IDA2 & 36,224.25 & 31,415.90 & 32.108 & 2.686 & 52.59 & 92.91\\
DA+IDA3 & 31,108.49 & 27,580.21 & 30.688 & 4.949 & 41.40 & 92.91\\
IDA1+IDA2 & 33,700.14 & 28,934.16 & 32.595 & 3.631 & 48.18 & 92.91\\
IDA1+IDA3 & 31,246.00 & 27,144.25 & 26.845 & 4.852 & 35.82 & 92.91\\
IDA2+IDA3 & 30,165.18 & 25,803.31 & 23.804 & 5.321 & 36.85 & 90.55\\
DA+IDA1+IDA2 & 42,720.38 & 37,787.34 & 32.665 & 3.132 & 52.59 & 92.91\\
DA+IDA1+IDA3 & 42,455.20 & 38,338.63 & 25.408 & 4.669 & 41.40 & 92.91\\
DA+IDA2+IDA3 & 45,935.23 & 41,734.28 & 23.345 & 4.315 & 36.85 & 89.78\\
IDA1+IDA2+IDA3 & 43,795.80 & 39,536.97 & 23.229 & 4.339 & 36.85 & 89.78\\
DA+IDA1+IDA2+IDA3 & 53,708.73 & 49,397.35 & 23.099 & 4.320 & 36.85 & 89.78\\
Sequential & 16,904.75 & 12,761.53 & 25.179 & 5.509 & 35.12 & 89.78\\
Perfect PV & 56,170.13 & 56,170.13 & 0.000 & 0.000 & 50.00 & 50.00\\
\bottomrule
\end{longtable}
\endgroup

\subsection{Farm 910}
\begingroup\small
\begin{longtable}{lrrrrrr}
\toprule
Profile & Profit D & Profit E & Abs. D & Abs. E & SOC D & SOC E\\
\midrule\endhead
DA & 19,804.62 & 16,479.18 & 85.420 & 27.383 & 13.26 & 7.00\\
IDA1 & 20,307.97 & 16,539.92 & 80.770 & 22.410 & 34.26 & 38.90\\
IDA2 & 20,814.37 & 17,130.78 & 79.143 & 19.284 & 50.00 & 61.69\\
IDA3 & 16,166.34 & 13,807.45 & 49.643 & 11.613 & 37.53 & 54.32\\
DA+IDA1 & 42,888.93 & 39,331.58 & 77.013 & 19.185 & 34.11 & 39.36\\
DA+IDA2 & 50,929.20 & 47,401.79 & 73.608 & 16.857 & 24.56 & 59.83\\
DA+IDA3 & 43,657.78 & 40,796.42 & 75.612 & 30.074 & 37.81 & 53.97\\
IDA1+IDA2 & 47,142.12 & 43,704.32 & 75.478 & 18.861 & 22.51 & 61.24\\
IDA1+IDA3 & 43,370.40 & 40,379.00 & 69.401 & 21.415 & 38.24 & 53.79\\
IDA2+IDA3 & 41,507.26 & 38,501.15 & 71.408 & 22.878 & 37.31 & 54.28\\
DA+IDA1+IDA2 & 61,319.41 & 57,815.61 & 76.593 & 18.861 & 28.59 & 60.53\\
DA+IDA1+IDA3 & 61,252.43 & 58,336.44 & 69.447 & 20.960 & 37.74 & 54.13\\
DA+IDA2+IDA3 & 66,905.52 & 63,917.51 & 70.746 & 20.507 & 37.69 & 54.32\\
IDA1+IDA2+IDA3 & 63,345.33 & 60,500.17 & 70.554 & 21.361 & 37.64 & 54.32\\
DA+IDA1+IDA2+IDA3 & 78,971.91 & 76,171.83 & 70.412 & 21.643 & 37.64 & 54.32\\
Sequential & 20,500.95 & 17,612.81 & 75.749 & 24.531 & 37.51 & 54.32\\
Perfect PV & 83,258.82 & 83,258.82 & 0.000 & 0.000 & 50.00 & 50.00\\
\bottomrule
\end{longtable}
\endgroup


\section{Earlier auction sensitivities and battery counterfactuals}
These are separate historical comparisons, not revisions to the matched
monthly table. Farm 1, January 2026, used the earlier physical/scenario setting.
The scalar penalty is a planning deterrent; the table reports actual cash
after signed settlement and wear, without deducting that illustrative penalty.
\begin{center}\small
\begin{tabular}{llrrr}
\toprule Formulation & PV information & Profit EUR & Abs. MWh & Wear EUR\\\midrule
Exact delivery & Forecast ensemble & 9,856.79 & 1.370 & 369.46\\
Exact delivery & Perfect PV & 10,582.72 & 0.000 & 224.23\\
Penalty 500 EUR/MWh & Forecast ensemble & 10,167.58 & 4.729 & 457.24\\
Penalty 500 EUR/MWh & Perfect PV & 10,582.72 & 0.000 & 224.23\\
Penalty 100 EUR/MWh & Forecast ensemble & 7,043.63 & 643.228 & 114.21\\
Penalty 100 EUR/MWh & Perfect PV & 7,112.42 & 642.110 & 90.22\\
\bottomrule\end{tabular}\end{center}
At the lower penalty the model takes very large short positions: forecast-case
auction receipts EUR 97,335.24 are offset by settlement EUR $-90,177.39$.
This is a strong objective-sensitivity diagnostic, not evidence that a
100 EUR/MWh coefficient represents actual imbalance economics. Even perfect
PV does not prevent deliberate planned deviations in that formulation.

In the separately saved January battery-minus-PV-only comparison, incremental
cash profit was EUR 57.77 (deterministic), 162.22 (51 paths), and 83.67 (perfect
PV). Added wear was EUR 520.36, 369.46 and 224.23; absolute imbalance changed
by $-5.041$, $-1.884$ and effectively zero MWh. These are operational cash
differences under the saved initial/terminal conditions, not annual battery
returns or capital-cost benefits.

The accompanying \path{auction_sensitivities.csv} retains 220 rows from these
comparisons and the farm-specific 10/20-day four-policy participation studies,
including ensemble-mean control, PV-only references, ending energy and saved
validation-warning counts where provided. Periods and objectives must be
matched before comparing rows. Sources added during the final audit carry
their own addition timestamp in the evidence ledger; the ongoing-search status
snapshot remains frozen at the original capture time.

\section{January MPC results by farm and case}
\label{app:mpc-farms}
Every number is target controller minus its own tracking baseline over the two
days January 15--16, 2026. End MWh is final stored-energy difference. These
are cash comparisons without inventory valuation. The CSV also preserves
absolute target/baseline totals and solver diagnostics.
\begingroup\small
\begin{longtable}{lrrrr}
\toprule
Farm / case & Profit EUR & Abs. MWh & Wear EUR & End MWh\\
\midrule\endhead
1 / adaptive\_continuation & 28.60 & 2.621 & 35.83 & 0.000\\
1 / adaptive\_ecc & 12.91 & 2.428 & 32.22 & 0.000\\
1 / adaptive\_monetary\_risk & -0.34 & 1.840 & 29.93 & 0.000\\
1 / adaptive\_raw & 18.40 & 2.389 & 31.49 & 0.000\\
1 / adaptive\_reduced & 14.57 & 1.730 & 19.06 & 0.000\\
1 / adaptive\_shuffled & 14.34 & 2.154 & 27.05 & 0.000\\
1 / calibrated\_double\_energy & 37.64 & 2.631 & 21.67 & 0.000\\
1 / calibrated\_half\_poi & 14.14 & 2.281 & 28.96 & 0.000\\
1 / calibrated\_half\_power & 20.10 & 1.337 & 10.35 & 0.000\\
1 / calibrated\_shuffled & 18.63 & 2.404 & 32.15 & 0.000\\
1 / causal\_auctions & 27.72 & 2.164 & 38.36 & 0.000\\
1 / causal\_fundamentals & 30.81 & 2.424 & 52.14 & 0.000\\
1 / da\_only\_frozen & 19.34 & 1.799 & 16.09 & 0.000\\
1 / da\_only\_latest & 10.17 & 2.218 & 24.87 & 0.000\\
1 / deterministic\_economic & 2.52 & 1.059 & -8.27 & 0.000\\
1 / deterministic\_oracle\_gap\_fallback & 27.51 & 0.387 & -0.58 & 0.000\\
1 / deterministic\_penalty & 25.48 & 0.412 & -0.08 & 0.000\\
1 / deterministic\_price\_oracle & 27.51 & 0.387 & -0.58 & 0.000\\
1 / double\_battery\_energy & 50.79 & 2.941 & 27.47 & 0.000\\
1 / ensemble\_economic & 13.24 & 2.273 & 29.31 & 0.000\\
1 / ensemble\_oracle\_gap\_fallback & 17.66 & 0.530 & 14.07 & 0.000\\
1 / ensemble\_penalty & 17.93 & 0.471 & 9.91 & 0.000\\
1 / ensemble\_price\_oracle & 17.66 & 0.530 & 14.07 & 0.000\\
1 / frozen\_information & 18.50 & 1.824 & 23.38 & 0.000\\
1 / half\_battery\_power & 26.68 & 1.492 & 13.24 & 0.000\\
1 / half\_poi\_capacity & 26.73 & 2.377 & 29.60 & 0.000\\
345 / adaptive\_continuation & 45.32 & 4.378 & -9.13 & -0.122\\
345 / adaptive\_ecc & 34.12 & 3.927 & -11.78 & -0.122\\
345 / adaptive\_monetary\_risk & 24.50 & 3.518 & -16.57 & -0.122\\
345 / adaptive\_raw & 34.32 & 3.871 & -13.77 & -0.122\\
345 / adaptive\_reduced & 46.06 & 3.863 & -12.05 & -0.122\\
345 / adaptive\_shuffled & 54.27 & 4.204 & -5.49 & -0.122\\
345 / calibrated\_double\_energy & 53.79 & 4.347 & -11.17 & -0.122\\
345 / calibrated\_half\_poi & 34.12 & 3.927 & -11.78 & -0.122\\
345 / calibrated\_half\_power & 50.26 & 2.748 & -30.57 & -0.122\\
345 / calibrated\_shuffled & 54.27 & 4.204 & -5.49 & -0.122\\
345 / causal\_auctions & 82.58 & 4.097 & 2.10 & -0.122\\
345 / causal\_fundamentals & 82.64 & 4.660 & 27.31 & -0.122\\
345 / da\_only\_frozen & 121.76 & 3.363 & -7.10 & -0.505\\
345 / da\_only\_latest & 117.86 & 3.910 & 5.49 & -0.505\\
345 / deterministic\_economic & 21.86 & 2.365 & -55.54 & -0.122\\
345 / deterministic\_oracle\_gap\_fallback & 21.39 & 1.293 & -27.49 & -0.028\\
345 / deterministic\_penalty & 16.40 & 1.199 & -27.49 & -0.028\\
345 / deterministic\_price\_oracle & 21.39 & 1.293 & -27.49 & -0.028\\
345 / double\_battery\_energy & 53.61 & 4.347 & -11.17 & -0.122\\
345 / ensemble\_economic & 30.46 & 3.333 & -23.91 & -0.122\\
345 / ensemble\_oracle\_gap\_fallback & 22.17 & 1.295 & -14.87 & -0.122\\
345 / ensemble\_penalty & 9.79 & 1.142 & -22.41 & -0.122\\
345 / ensemble\_price\_oracle & 22.17 & 1.295 & -14.87 & -0.122\\
345 / frozen\_information & 47.78 & 3.386 & -19.88 & -0.098\\
345 / half\_battery\_power & 50.17 & 2.748 & -30.57 & -0.122\\
345 / half\_poi\_capacity & 34.39 & 3.837 & -14.41 & -0.122\\
901 / adaptive\_continuation & 74.79 & 5.413 & 62.53 & 0.000\\
901 / adaptive\_ecc & 70.93 & 5.051 & 64.15 & 0.000\\
901 / adaptive\_monetary\_risk & 8.45 & 4.700 & 81.96 & 0.000\\
901 / adaptive\_raw & 76.12 & 4.733 & 58.32 & 0.000\\
901 / adaptive\_reduced & -50.15 & 6.956 & 120.50 & 0.000\\
901 / adaptive\_shuffled & 80.62 & 6.135 & 98.63 & 0.000\\
901 / calibrated\_double\_energy & 119.15 & 6.586 & 73.97 & 0.000\\
901 / calibrated\_half\_poi & 64.59 & 5.030 & 66.06 & 0.000\\
901 / calibrated\_half\_power & 72.77 & 3.419 & 29.90 & 0.000\\
901 / calibrated\_shuffled & 78.74 & 5.717 & 86.96 & 0.000\\
901 / causal\_auctions & 77.48 & 4.564 & 61.47 & 0.000\\
901 / causal\_fundamentals & 29.66 & 6.093 & 117.38 & 0.000\\
901 / da\_only\_frozen & 46.09 & 4.008 & 26.41 & 0.000\\
901 / da\_only\_latest & 67.54 & 4.538 & 33.84 & 0.000\\
901 / deterministic\_economic & 19.22 & 2.426 & -14.61 & 0.000\\
901 / deterministic\_oracle\_gap\_fallback & 5.22 & 1.502 & 32.07 & 0.365\\
901 / deterministic\_penalty & -7.48 & 1.653 & 32.07 & 0.365\\
901 / deterministic\_price\_oracle & 5.22 & 1.502 & 32.07 & 0.365\\
901 / double\_battery\_energy & 91.73 & 6.579 & 75.43 & 0.000\\
901 / ensemble\_economic & 36.38 & 4.236 & 48.64 & 0.000\\
901 / ensemble\_oracle\_gap\_fallback & -140.68 & 3.270 & 142.36 & 0.000\\
901 / ensemble\_penalty & -68.02 & 2.818 & 77.99 & 0.000\\
901 / ensemble\_price\_oracle & -140.68 & 3.270 & 142.36 & 0.000\\
901 / frozen\_information & 44.34 & 4.296 & 40.17 & 0.000\\
901 / half\_battery\_power & 73.06 & 3.458 & 30.63 & 0.000\\
901 / half\_poi\_capacity & 76.12 & 4.733 & 58.32 & 0.000\\
910 / adaptive\_continuation & 182.75 & 15.383 & 208.73 & 0.000\\
910 / adaptive\_ecc & 184.10 & 12.994 & 172.20 & 0.000\\
910 / adaptive\_monetary\_risk & 62.23 & 11.019 & 189.31 & 0.000\\
910 / adaptive\_raw & 172.89 & 12.557 & 180.47 & 0.000\\
910 / adaptive\_reduced & -140.79 & 19.639 & 437.05 & 0.000\\
910 / adaptive\_shuffled & 183.04 & 13.995 & 225.56 & 0.000\\
910 / calibrated\_double\_energy & 286.48 & 16.529 & 205.19 & 0.000\\
910 / calibrated\_half\_poi & 209.10 & 13.419 & 172.20 & 0.000\\
910 / calibrated\_half\_power & 147.01 & 8.000 & 108.56 & 0.000\\
910 / calibrated\_shuffled & 204.34 & 14.659 & 229.21 & 0.000\\
910 / causal\_auctions & 112.76 & 10.840 & 201.52 & 0.000\\
910 / causal\_fundamentals & 106.27 & 14.371 & 345.59 & 0.000\\
910 / da\_only\_frozen & 95.52 & 10.855 & 75.03 & 0.000\\
910 / da\_only\_latest & 65.00 & 12.275 & 156.49 & 0.000\\
910 / deterministic\_economic & 54.50 & 6.444 & -24.70 & 0.000\\
910 / deterministic\_oracle\_gap\_fallback & 164.31 & 2.391 & 12.64 & 0.026\\
910 / deterministic\_penalty & 89.27 & 2.442 & 13.33 & 0.000\\
910 / deterministic\_price\_oracle & 164.31 & 2.391 & 12.64 & 0.026\\
910 / double\_battery\_energy & 282.05 & 15.904 & 205.19 & 0.000\\
910 / ensemble\_economic & 97.73 & 10.942 & 133.42 & 0.000\\
910 / ensemble\_oracle\_gap\_fallback & 109.43 & 5.469 & 165.74 & 0.000\\
910 / ensemble\_penalty & 115.11 & 4.682 & 120.60 & 0.000\\
910 / ensemble\_price\_oracle & 109.43 & 5.469 & 165.74 & 0.000\\
910 / frozen\_information & 122.86 & 12.038 & 111.99 & 0.000\\
910 / half\_battery\_power & 144.80 & 7.687 & 108.56 & 0.000\\
910 / half\_poi\_capacity & 172.61 & 12.622 & 172.20 & 0.000\\
\bottomrule
\end{longtable}
\endgroup


\section{Price experiment inventory}
\label{app:price-inventory}
Identifiers below are relative to the MPC tree's
\path{outputs/energy_price_forecasts/}. ``Records'' counts manifest result/model
entries, not independent observations or total weekly refits. Derived comparison
reports can have zero such entries while retaining completed group metrics.
Status is frozen at capture. Complete does not certify numerical convergence,
causal availability or untouched evaluation. Per-run periods and all candidate
metrics are in \path{price_run_inventory.csv} and
\path{price_candidate_results.csv} beside the snapshot.
\begingroup\small
\begin{longtable}{p{0.71\textwidth}lr}
\toprule
Saved run identifier & Status & Records\\
\midrule\endhead
\path{chronological_2026_v3} & complete & 30\\
\path{european_context_ablation_v1} & complete & 75\\
\path{european_context_previous_day_v1} & complete & 75\\
\path{extended_2024_training_v2} & complete & 27\\
\path{nord_da_2026_huber_search_v1} & complete & 56\\
\path{nord_da_2026_previous_ida_prices_v1} & complete & 140\\
\path{nord_da_2026_previous_idas_ablation_v1} & complete & 0\\
\path{nord_da_2026_previous_idas_selected_rolling_v1} & complete & 10\\
\path{nord_da_2026_recent_training_windows_v1} & complete & 96\\
\path{nord_da_2026_recent_windows_selected_rolling_v1} & running & 5\\
\path{nord_da_curve_monthly_diagnostics_v1} & complete & 0\\
\path{nord_da_curves_rolling_v1} & complete & 15\\
\path{nord_da_daily_curves_search_v1} & complete & 315\\
\path{nord_da_expanded_context_search_v1} & complete & 243\\
\path{nord_da_french_balance_matched_v1} & complete & 108\\
\path{nord_da_per_slot_search_v1} & complete & 40\\
\path{nord_da_resolution_holdout_v1} & complete & 10\\
\path{nord_da_resolution_search_v1} & complete & 60\\
\path{nord_ida_2024_hourly_exports_v1} & complete & 243\\
\path{nord_ida_2025_selected_rolling_holdout_v1} & complete & 27\\
\path{nord_ida_2025_terna_load_ablation_v1} & complete & 0\\
\path{nord_ida_2025_terna_load_matched_v1} & complete & 324\\
\path{nord_ida_2025_terna_selected_rolling_holdout_v1} & complete & 21\\
\path{nord_ida_2026_huber_search_v1} & complete & 144\\
\path{nord_ida_bidirectional_neighbors_search_v1} & running & 470\\
\path{nord_ida_cleared_context_2025_v1} & complete & 486\\
\path{nord_ida_cleared_context_search_v1} & complete & 486\\
\path{nord_ida_cleared_monthly_diagnostics_v1} & complete & 0\\
\path{nord_ida_cleared_rolling_v1} & complete & 27\\
\path{nord_ida_french_balance_ablation_v1} & complete & 0\\
\path{nord_ida_french_balance_matched_v1} & complete & 324\\
\path{nord_ida_nearest_auction_spreads_v1} & complete & 324\\
\path{nord_ida_per_slot_search_v1} & complete & 72\\
\path{nord_ida_swiss_border_matched_v1} & complete & 324\\
\path{nord_ida_swiss_reservoir_ablation_v1} & complete & 0\\
\path{nord_ida_swiss_reservoir_matched_v1} & complete & 324\\
\path{nord_ida_temporal_ensemble_v1} & complete & 0\\
\path{nord_ida_temporal_hourwise_v1} & complete & 0\\
\path{nord_ida_temporal_stacking_v1} & complete & 0\\
\path{nord_intraday_settlement_resolution_search_v1} & complete & 120\\
\bottomrule
\end{longtable}
\endgroup


\section{Evidence ledger and review method}
The report uses the full main LaTeX as its base, reconciles the sibling
feature LaTeX and MPC/price LaTeX, and checks the optimizer, physical helpers,
information ownership, replay accounting, forecast adapters and effective
configuration. Numerical evidence is read from completed run records, compact
CSV exports and saved independent audits; no optimization or training was
rerun for this report. Where only an older report survives, that weaker evidence
is stated. The source ledger below maps every pinned input to its checkout;
full SHA-256 values and captured working-tree state are in
\path{docs/report_evidence/2026-11-04/snapshot.json}.

The frozen CSV tables are a compact report snapshot, not replacement live run
manifests. They preserve comparison labels and target provenance. The main
working tree had unrelated ongoing PV-training modifications and an existing
price-data inventory; those are outside this documentation change. The two
sibling working trees were read as evidence rather than modified or merged.
The independently written presentation document is
\path{docs/research_overview.tex}; it retains the scientific assumptions without
implementation identifiers and places its detailed formulation in appendices.
\begingroup\small
\begin{longtable}{p{0.16\textwidth}p{0.76\textwidth}}
\toprule ID / checkout & Evidence path\\\midrule\endhead
S001 / main & \path{docs/results/pv_reviewed_hpo/historical_ghi_all_farms_night3000_v1/top_three_metrics.csv}\\
S002 / main & \path{outputs/continuous_optimization/comparisons/farm_1/2026_01_mixed_models_pv_comparison/monthly_comparison.csv}\\
S003 / main & \path{outputs/continuous_optimization/comparisons/farm_1/2026_08_mixed_models_pv_comparison/monthly_comparison.csv}\\
S004 / main & \path{outputs/continuous_optimization/comparisons/farm_345/2026_01_mixed_models_pv_comparison/monthly_comparison.csv}\\
S005 / main & \path{outputs/continuous_optimization/comparisons/farm_345/2026_08_mixed_models_pv_comparison/monthly_comparison.csv}\\
S006 / main & \path{outputs/continuous_optimization/comparisons/farm_901/2026_01_mixed_models_pv_comparison/monthly_comparison.csv}\\
S007 / main & \path{outputs/continuous_optimization/comparisons/farm_901/2026_08_mixed_models_pv_comparison/monthly_comparison.csv}\\
S008 / main & \path{outputs/continuous_optimization/comparisons/farm_910/2026_01_mixed_models_pv_comparison/monthly_comparison.csv}\\
S009 / main & \path{outputs/continuous_optimization/comparisons/farm_910/2026_08_mixed_models_pv_comparison/monthly_comparison.csv}\\
S010 / main & \path{outputs/continuous_optimization/comparisons/farm_1/2026_01_continuous_participation_forecast_comparison/monthly_comparison.csv}\\
S011 / main & \path{outputs/continuous_optimization/comparisons/farm_345/2026_01_continuous_participation_forecast_comparison/monthly_comparison.csv}\\
S012 / main & \path{outputs/continuous_optimization/comparisons/farm_901/2026_01_continuous_participation_forecast_comparison/monthly_comparison.csv}\\
S013 / main & \path{outputs/continuous_optimization/comparisons/farm_910/2026_01_continuous_participation_forecast_comparison/monthly_comparison.csv}\\
S014 / MPC & \path{outputs/energy_research/reviewed_january_matrix_v1/manifest.json}\\
S015 / MPC & \path{outputs/energy_price_forecasts/chronological_2026_v3/manifest.json}\\
S016 / MPC & \path{outputs/energy_price_forecasts/european_context_ablation_v1/manifest.json}\\
S017 / MPC & \path{outputs/energy_price_forecasts/european_context_previous_day_v1/manifest.json}\\
S018 / MPC & \path{outputs/energy_price_forecasts/extended_2024_training_v2/manifest.json}\\
S019 / MPC & \path{outputs/energy_price_forecasts/nord_da_2026_huber_search_v1/manifest.json}\\
S020 / MPC & \path{outputs/energy_price_forecasts/nord_da_2026_previous_ida_prices_v1/manifest.json}\\
S021 / MPC & \path{outputs/energy_price_forecasts/nord_da_2026_previous_idas_ablation_v1/manifest.json}\\
S022 / MPC & \path{outputs/energy_price_forecasts/nord_da_2026_previous_idas_selected_rolling_v1/manifest.json}\\
S023 / MPC & \path{outputs/energy_price_forecasts/nord_da_2026_recent_training_windows_v1/manifest.json}\\
S024 / MPC & \path{outputs/energy_price_forecasts/nord_da_2026_recent_windows_selected_rolling_v1/manifest.json}\\
S025 / MPC & \path{outputs/energy_price_forecasts/nord_da_curve_monthly_diagnostics_v1/manifest.json}\\
S026 / MPC & \path{outputs/energy_price_forecasts/nord_da_curves_rolling_v1/manifest.json}\\
S027 / MPC & \path{outputs/energy_price_forecasts/nord_da_daily_curves_search_v1/manifest.json}\\
S028 / MPC & \path{outputs/energy_price_forecasts/nord_da_expanded_context_search_v1/manifest.json}\\
S029 / MPC & \path{outputs/energy_price_forecasts/nord_da_french_balance_matched_v1/manifest.json}\\
S030 / MPC & \path{outputs/energy_price_forecasts/nord_da_per_slot_search_v1/manifest.json}\\
S031 / MPC & \path{outputs/energy_price_forecasts/nord_da_resolution_holdout_v1/manifest.json}\\
S032 / MPC & \path{outputs/energy_price_forecasts/nord_da_resolution_search_v1/manifest.json}\\
S033 / MPC & \path{outputs/energy_price_forecasts/nord_ida_2024_hourly_exports_v1/manifest.json}\\
S034 / MPC & \path{outputs/energy_price_forecasts/nord_ida_2025_selected_rolling_holdout_v1/manifest.json}\\
S035 / MPC & \path{outputs/energy_price_forecasts/nord_ida_2025_terna_load_ablation_v1/manifest.json}\\
S036 / MPC & \path{outputs/energy_price_forecasts/nord_ida_2025_terna_load_matched_v1/manifest.json}\\
S037 / MPC & \path{outputs/energy_price_forecasts/nord_ida_2025_terna_selected_rolling_holdout_v1/manifest.json}\\
S038 / MPC & \path{outputs/energy_price_forecasts/nord_ida_2026_huber_search_v1/manifest.json}\\
S039 / MPC & \path{outputs/energy_price_forecasts/nord_ida_bidirectional_neighbors_search_v1/manifest.json}\\
S040 / MPC & \path{outputs/energy_price_forecasts/nord_ida_cleared_context_2025_v1/manifest.json}\\
S041 / MPC & \path{outputs/energy_price_forecasts/nord_ida_cleared_context_search_v1/manifest.json}\\
S042 / MPC & \path{outputs/energy_price_forecasts/nord_ida_cleared_monthly_diagnostics_v1/manifest.json}\\
S043 / MPC & \path{outputs/energy_price_forecasts/nord_ida_cleared_rolling_v1/manifest.json}\\
S044 / MPC & \path{outputs/energy_price_forecasts/nord_ida_french_balance_ablation_v1/manifest.json}\\
S045 / MPC & \path{outputs/energy_price_forecasts/nord_ida_french_balance_matched_v1/manifest.json}\\
S046 / MPC & \path{outputs/energy_price_forecasts/nord_ida_nearest_auction_spreads_v1/manifest.json}\\
S047 / MPC & \path{outputs/energy_price_forecasts/nord_ida_per_slot_search_v1/manifest.json}\\
S048 / MPC & \path{outputs/energy_price_forecasts/nord_ida_swiss_border_matched_v1/manifest.json}\\
S049 / MPC & \path{outputs/energy_price_forecasts/nord_ida_swiss_reservoir_ablation_v1/manifest.json}\\
S050 / MPC & \path{outputs/energy_price_forecasts/nord_ida_swiss_reservoir_matched_v1/manifest.json}\\
S051 / MPC & \path{outputs/energy_price_forecasts/nord_ida_temporal_ensemble_v1/manifest.json}\\
S052 / MPC & \path{outputs/energy_price_forecasts/nord_ida_temporal_hourwise_v1/manifest.json}\\
S053 / MPC & \path{outputs/energy_price_forecasts/nord_ida_temporal_stacking_v1/manifest.json}\\
S054 / MPC & \path{outputs/energy_price_forecasts/nord_intraday_settlement_resolution_search_v1/manifest.json}\\
S055 / features & \path{outputs/price_benchmark/20261023T133005867139Z/results.csv}\\
S056 / main & \path{outputs/pv_reviewed_hpo/four_cycle_neural_night3000_gpu_parallel_v1/status.json}\\
S057 / main & \path{outputs/pv_reviewed_hpo/four_cycle_neural_night3000_gpu_v1/status.json}\\
S058 / main & \path{outputs/pv_reviewed_hpo/four_cycle_neural_night3000_multi_gpu_v1/status.json}\\
S059 / main & \path{outputs/pv_reviewed_hpo/four_cycle_neural_night3000_v1/status.json}\\
S060 / main & \path{outputs/pv_reviewed_hpo/four_cycle_neural_night3000_v2/status.json}\\
S061 / main & \path{outputs/pv_reviewed_hpo/historical_ghi_all_farms_night3000_v1/status.json}\\
S062 / main & \path{outputs/pv_reviewed_hpo/historical_ghi_all_farms_v1/status.json}\\
S063 / main & \path{AGENTS.md}\\
S064 / main & \path{configs/standard.yaml}\\
S065 / main & \path{docs/pv_probabilistic_online_crps_reference.tex}\\
S066 / main & \path{docs/results/pv_online_crps_january_august_2026/metrics.csv}\\
S067 / main & \path{docs/results/probabilistic_pv_january_august_2026/comparison/paired_comparison.csv}\\
S068 / main & \path{docs/results/pv_realized_irradiance_2025_2026/README.md}\\
S069 / main & \path{docs/results/cleaned_pv_hpo/all_farms_cleaned_16_trials/selected_model_summary.csv}\\
S070 / main & \path{docs/results/realtime_mpc_2026_01_15_16/continuous_controller_comparison.csv}\\
S071 / main & \path{docs/results/realtime_mpc_2026_01_15_16/frozen_commitments_controller_comparison.csv}\\
S072 / main & \path{outputs/realtime_mpc/may_2026_all_farms/MAY_2026_ALL_FARMS_SUMMARY.md}\\
S073 / main & \path{models/pv/reviewed_ghi_2026_selected/manifest.json}\\
S074 / main & \path{outputs/continuous_optimization/comparisons/farm_1/envelope_replay_2026_08_02_to_11/period_comparison.csv}\\
S075 / main & \path{outputs/continuous_optimization/comparisons/farm_1/envelope_replay_2026_08_02_to_11/matched_state_solves.csv}\\
S076 / main & \path{outputs/participation_studies/gme_nord_farm910_2025-01_current_settings_sequential_ida3_noon/json/manifest.json}\\
S077 / main & \path{outputs/participation_studies/gme_nord_farm910_2025-01_current_settings_sequential_ida3_noon/csv/profile_ranking.csv}\\
S078 / features & \path{docs/mathematical_model.tex}\\
S079 / features & \path{docs/price_benchmark.md}\\
S080 / features & \path{outputs/price_features/history_2024_coverage.json}\\
S081 / features & \path{outputs/price_benchmark/20261023T133005867139Z/manifest.json}\\
S082 / features & \path{outputs/price_benchmark/20261023T133151054322Z/results.csv}\\
S083 / features & \path{outputs/price_benchmark/20261023T134106810433Z/results.csv}\\
S084 / features & \path{outputs/price_benchmark/20261023T134324360680Z/results.csv}\\
S085 / MPC & \path{docs/realtime_mpc_price_forecasting.tex}\\
S086 / MPC & \path{docs/energy_market_research_extensions.md}\\
S087 / MPC & \path{docs/reviews/mpc_results_overview.md}\\
S088 / MPC & \path{docs/reviews/price_huber_convergence.md}\\
S089 / MPC & \path{docs/reviews/price_da_2026_previous_idas.md}\\
S090 / MPC & \path{docs/reviews/price_ida_2025_selected_rolling.md}\\
S091 / MPC & \path{docs/reviews/continuous_mpc_reviewed_october_v1.json}\\
S092 / MPC & \path{docs/reviews/continuous_mpc_reviewed_farm910_august_v1.json}\\
S093 / MPC & \path{outputs/energy_research/reviewed_january_matrix_audit.json}\\
S094 / MPC & \path{outputs/research_inputs/pv_calibration_reviewed/calibration_summary.json}\\
S095 / main & \path{src/optimization/model.py}\\
S096 / main & \path{src/optimization/mpc.py}\\
S097 / main & \path{src/optimization/physical.py}\\
S098 / main & \path{src/optimization/degradation.py}\\
S099 / main & \path{src/optimization/pv_envelope.py}\\
S100 / main & \path{src/simulation/replay.py}\\
S101 / main & \path{src/simulation/continuous.py}\\
S102 / main & \path{src/simulation/continuous_inputs.py}\\
S103 / main & \path{src/forecasting/inference/pv_reviewed_hpo.py}\\
S104 / features & \path{src/forecasting/price_benchmark/models.py}\\
S105 / features & \path{src/forecasting/price_benchmark/data.py}\\
S106 / features & \path{src/forecasting/price_benchmark/evaluation.py}\\
S107 / MPC & \path{src/optimization/mpc.py}\\
S108 / MPC & \path{src/optimization/adaptive_mpc.py}\\
S109 / MPC & \path{src/scenario_research.py}\\
S110 / MPC & \path{src/forecasting/continuation.py}\\
S111 / MPC & \path{src/forecasting/energy_prices.py}\\
S112 / MPC & \path{src/forecasting/price_search.py}\\
S113 / MPC & \path{src/forecasting/price_search_models.py}\\
S114 / main & \path{outputs/continuous_optimization/comparisons/farm_1/2026_01_participation_days_20_four_policies/period_comparison.csv}\\
S115 / main & \path{outputs/continuous_optimization/comparisons/farm_1/2026_01_participation_ten_days_four_policies/period_comparison.csv}\\
S116 / main & \path{outputs/continuous_optimization/comparisons/farm_1/2026_01_pv_battery_comparison/monthly_comparison.csv}\\
S117 / main & \path{outputs/continuous_optimization/comparisons/farm_1/2026_01_pv_deterministic_vs_scenarios/monthly_comparison.csv}\\
S118 / main & \path{outputs/continuous_optimization/comparisons/farm_1/2026_01_pv_scenarios_comparison/monthly_comparison.csv}\\
S119 / main & \path{outputs/continuous_optimization/comparisons/farm_1/2026_01_pv_scenarios_penalty100_comparison/monthly_comparison.csv}\\
S120 / main & \path{outputs/continuous_optimization/comparisons/farm_345/2026_01_participation_ten_days_four_policies/period_comparison.csv}\\
S121 / main & \path{outputs/continuous_optimization/comparisons/farm_901/2026_01_participation_ten_days_four_policies/period_comparison.csv}\\
S122 / main & \path{outputs/continuous_optimization/comparisons/farm_910/2026_01_participation_ten_days_four_policies/period_comparison.csv}\\
S123 / main & \path{outputs/continuous_optimization/comparisons/farm_1/2026_01_pv_battery_comparison/battery_effect_by_policy.csv}\\
S124 / MPC & \path{src/forecasting/pv_scenario_calibration.py}\\
\bottomrule\end{longtable}\endgroup

